\section{Daten, Vorverarbeitung und Agentendesign}
\label{sec:agenten_design}

\subsection{Datenquelle und Trainingsdaten}

Anders als bei überwachten Lernverfahren existiert kein vorab
erhobener Datensatz. Sämtliche Übergänge $(s_t,a_t,r_{t+1},s_{t+1})$
entstehen online durch die Interaktion zwischen Agent und selbst
implementierter Simulationsumgebung. Die Trainingsdatenverteilung
hängt damit von der jeweils aktuellen $\epsilon$-greedy Policy ab. Zu
Beginn erzeugt die hohe Exploration vielfältige Übergänge; im
späteren Training dominiert zunehmend die gelernte Policy.

Die Verwendung einer Simulation erlaubt eine sehr große Zahl
kostengünstiger Episoden und vollständig bekannte Rewards. Sie
erzeugt zugleich ein Validitätsrisiko: Der Agent optimiert exakt die
implementierten Regeln und kann Modellvereinfachungen ausnutzen, die
in einem realen Blackjack-Spiel nicht existieren.

\subsection{Zustandsrepräsentation und Vorverarbeitung}

Die beiden Agenten unterscheiden sich ausschließlich durch ihren
State Encoder. Dadurch lässt sich der Nutzen der zusätzlichen
Kontextmerkmale isolierter untersuchen als bei gleichzeitig
verschiedenen Lernalgorithmen.

\paragraph{Baseline-Agent.}
Der Baseline-Zustand ist
\begin{equation}
  s_t^{B}=(p_t,d_t,a_t).
\end{equation}
Alle drei Werte werden als Ganzzahlen gespeichert. Die finale
Q-Tabelle enthält für jeden Trainingsseed genau 280 beobachtete Zustände.

\paragraph{Counting-Agent.}
Der erweiterte Zustand ist
\begin{equation}
  s_t^{C}=(p_t,d_t,a_t,\tilde r_t,\tilde c_t,\tilde m_t).
\end{equation}
Zur Begrenzung des Zustandsraums werden die Zusatzmerkmale
diskretisiert. Dabei bezeichnet $\operatorname{trunc}$ die im Code
verwendete Ganzzahlumwandlung, die Nachkommastellen gegen null abschneidet:
\begin{align}
  \tilde r_t &= \operatorname{trunc}(\operatorname{clip}(r_t,-20,20)),\\
  \tilde c_t &= \operatorname{trunc}(\operatorname{clip}(c_t,-10,10)),\\
  \tilde m_t &= \operatorname{clip}(\lfloor m_t/52\rfloor,0,6).
\end{align}
Eine klassische Normalisierung ist für die Q-Tabelle nicht
erforderlich, weil keine kontinuierliche Funktion approximiert wird.
Das Clipping verhindert seltene Extremwerte, führt aber zu
Informationsverlust. Zudem sind Running Count, True Count und
verbleibende Decks mathematisch voneinander abhängig. Die
Repräsentation enthält somit Redundanz und vergrößert die Tabelle stark.

Am Ende des finalen Trainings enthalten die drei Counting-Tabellen
zwischen 113.744 und 113.769 Zustände. Im Mittel sind dies 113.760
und damit rund 406-mal so viele wie bei der Baseline. Nicht besuchte
Zustände erhalten implizit Q-Werte von null. Bei Gleichstand wählt
\texttt{argmax} die Aktion Stand; diese technische Tie-Break-Regel
beeinflusst insbesondere unbekannte oder kaum trainierte Zustände.

\subsection{Modellarchitektur}

Es wird kein Deep-Learning-Modell eingesetzt. Die Architektur besteht
aus einer Python-\texttt{defaultdict}, die für jeden beobachteten
Zustand einen Vektor mit zwei Q-Werten speichert:
\begin{equation}
  Q(s)=[Q(s,\text{Stand}),Q(s,\text{Hit})].
\end{equation}
Neue Zustände werden mit $[0,0]$ initialisiert. Nach jedem Schritt
erfolgt genau ein Online-Update gemäß Gleichung~\ref{eq:qlearning}.
Es gibt weder Mini-Batches noch Experience Replay oder Target
Network. Der Vorteil dieser Architektur ist ihre Nachvollziehbarkeit;
ihr Nachteil ist die fehlende Generalisierung zwischen ähnlichen Zuständen.

\subsection{Hyperparameter und Auswahlprozess}

Vor dem finalen Lauf wurden kombinierte Hyperparameterkonfigurationen
in drei Stufen getestet. Die Parameter wurden nicht einzeln in einer
vollständigen Ablationsstudie variiert; die Ergebnisse zeigen daher
eine pragmatische Auswahl, aber kein globales Optimum.

\begin{table}[H]
  \centering
  \caption{Stufen der Hyperparametersuche}
  \label{tab:tuning_stages}
  \small
  \begin{tabularx}{\textwidth}{lrrrX}
    \toprule
    \textbf{Stufe} & \textbf{Konfigurationen} &
    \textbf{Trainingsseeds} & \textbf{Training} & \textbf{Greedy Evaluation}\\
    \midrule
    A & 16 & 1 & 3 Mio. Episoden & 150.000 je Checkpoint\\
    B & 8  & 2 & 10 Mio. Episoden & 500.000 je Checkpoint\\
    C & 4 Counting + Baseline & 3 & 20 Mio. Episoden & 1 Mio. je Checkpoint\\
    \bottomrule
  \end{tabularx}
\end{table}

Für den finalen Lauf wurde die Counting-Konfiguration
\texttt{c05\_long\_default} ausgewählt. Sie war in Stufe C mit einem
mittleren Reward von $-0{,}05259$ die beste Counting-Variante und
zeigte mit einer Seed-Standardabweichung von $0{,}00052$ die höchste
Stabilität der untersuchten Counting-Konfigurationen. Die Baseline
lag bei 20 Millionen Episoden mit $-0{,}04906$ dennoch vorn. Das
sprach dafür, dem größeren Counting-Zustandsraum erheblich mehr
Training zu geben.

\begin{table}[H]
  \centering
  \caption{Hyperparameter des finalen Trainings}
  \label{tab:hyperparameters}
  \begin{tabularx}{\textwidth}{lX}
    \toprule
    \textbf{Hyperparameter} & \textbf{Wert und Umsetzung}\\
    \midrule
    Lernrate $\alpha$ & $0{,}01$, während des gesamten Trainings konstant\\
    Diskontierungsfaktor $\gamma$ & $0{,}99$\\
    Initiales $\epsilon$ & $1{,}0$\\
    Finales $\epsilon$ & $0{,}10$\\
    $\epsilon$-Decay & Linear über 90\,\% der 100 Mio. Episoden,
    danach konstant\\
    Trainingsepisoden & 100 Mio. je Agentinstanz\\
    Checkpoint-Intervall & 10 Mio. Episoden\\
    Batch Size & Nicht anwendbar, da Online-Update nach jedem Schritt\\
    \bottomrule
  \end{tabularx}
\end{table}

Die finale Konfiguration wurde für beide Agententypen verwendet, um
den Modellvergleich nicht zusätzlich durch unterschiedliche
Lernparameter zu verzerren.
